\chapter{Derivatives of Exponentials, Logarithms, and Trig}

So far the derivative has been a machine that eats polynomials and rational
functions and spits out their slopes. The power rule, the product rule, the
quotient rule, the chain rule -- these handle everything you can build from
$x$ by adding, multiplying, and composing. But the most \emph{important}
functions in all of science are not built that way. Populations grow like
$e^{kt}$. Radioactive atoms decay like $e^{-\lambda t}$. A pendulum swings
like $\cos(\omega t)$. The loudness you hear and the brightness you see are
logarithmic. If calculus could not differentiate $e^x$, $\ln x$, $\sin x$,
and their relatives, it would be a beautiful theory about almost nothing.

This chapter computes the derivatives of all of them. But I want to do more
than hand you a table to memorize. Every one of these derivatives is
\emph{forced} by the definition of the function. The fact that
$\ddx{e^x} = e^x$ is not a lucky accident -- it is the whole reason the
number $e$ exists. The fact that $\ddx{\sin x} = \cos x$ is the unit circle
turning. By the end you should feel that these formulas could not have come
out any other way.

\chapterline

\section{The Exponential: a Function Equal to Its Own Slope}
\label{ch12:exp-intro}

Start with the most consequential single fact in differential
calculus. Take any exponential function $f(x) = b^x$ with base $b > 0$ and
ask for its derivative straight from the definition.

\begin{align*}
f'(x)
&= \lim_{h \to 0} \frac{b^{x+h} - b^x}{h} \\[2pt]
&= \lim_{h \to 0} \frac{b^x b^h - b^x}{h}
   \qquad\text{(law of exponents: } b^{x+h}=b^x b^h\text{)}\\[2pt]
&= \lim_{h \to 0}\; b^x \cdot \frac{b^h - 1}{h} \\[2pt]
&= b^x \cdot \underbrace{\lim_{h \to 0} \frac{b^h - 1}{h}}_{\textstyle =\, m(b)}.
\end{align*}

Stop and look at what just happened, because it is remarkable. The $b^x$
factored straight out of the limit, since it does not depend on $h$. What is
left over is a limit that does not depend on $x$ at all -- it is just a
\emph{number}, a constant that depends only on the base $b$. Call it $m(b)$.
So we have proven, with almost no work, the astonishing structural fact:
\[
   \ddx{b^x} \;=\; m(b)\cdot b^x,
   \qquad
   m(b) = \lim_{h \to 0}\frac{b^h - 1}{h}.
\]

\begin{focusbox}
The derivative of \emph{any} exponential function is a constant multiple of
\emph{itself}. The function $b^x$ reproduces its own shape when
differentiated; the only thing the derivative does is rescale it by the
fixed number $m(b)$. No other family of functions behaves like this, and it
is the source of the exponential's grip on the natural world.
\end{focusbox}

\begin{intuitionbox}
Why should the slope of $b^x$ be proportional to its height? Picture
compound interest. The more money you have, the more interest you earn per
year -- the growth is proportional to the current amount. The more rabbits
you have, the more baby rabbits per month. Anything whose rate of change is
proportional to its current size is, secretly, an exponential. The constant
$m(b)$ is just the proportionality factor, the ``interest rate'' baked into
the base.
\end{intuitionbox}

\subsection{What is that mysterious constant $m(b)$?}

The number $m(b) = \lim_{h\to 0}\frac{b^h-1}{h}$ has a clean geometric
meaning. Set $x = 0$ in our formula: since $b^0 = 1$,
\[
   f'(0) = m(b)\cdot b^0 = m(b).
\]
So $m(b)$ is nothing but \emph{the slope of $b^x$ at the point where it
crosses the $y$-axis}. Every exponential passes through $(0,1)$; what
distinguishes them is how steeply they leave that point. And our boxed
formula says something almost too good: the slope everywhere is the
$y$-axis slope times the height. Know the starting slope, and you know every
slope.

We can estimate $m(b)$ for a few bases by plugging a tiny $h$ into
$(b^h-1)/h$.

\begin{labbox}
Take $h = 0.001$ and compute $(b^h - 1)/h$ with a calculator.
\begin{itemize}
\item $b = 2$: \; $(2^{0.001}-1)/0.001 \approx 0.6934$.
\item $b = 2.5$: \; $(2.5^{0.001}-1)/0.001 \approx 0.9167$.
\item $b = 2.718$:\; $(2.718^{0.001}-1)/0.001 \approx 1.0004$.
\item $b = 3$: \; $(3^{0.001}-1)/0.001 \approx 1.0992$.
\end{itemize}
The slope $m(b)$ slides smoothly from below $1$ (at $b=2$) to above $1$
(at $b=3$), crossing $1$ somewhere in between. At $b=2.718$ the ratio is
already a whisker over $1$ -- so the true crossing point, where the
\emph{limit} $m(b)$ equals exactly $1$, is staring at us right around
$2.718$. (The finite step $h=0.001$ overshoots the limit slightly, which is
why $2.718$ reads $1.0004$ rather than just under $1$; shrink $h$ and the
numbers settle onto $m(b)=\ln b$, the value we will prove shortly.)
\end{labbox}

\begin{figure}[H]
\centering
\begin{tikzpicture}
\begin{axis}[
  width=11cm, height=7cm,
  axis lines=middle,
  xlabel=$x$, ylabel=$y$,
  xmin=-2.2, xmax=1.8, ymin=-0.3, ymax=6,
  domain=-2.2:1.8, samples=120,
  legend pos=north west, legend cell align=left,
  every axis plot/.append style={very thick},
]
\addplot[curveblue]   {2^x};   \addlegendentry{$2^x$ (slope $<1$ at $0$)}
\addplot[tangentred]  {2.71828^x}; \addlegendentry{$e^x$ (slope $=1$ at $0$)}
\addplot[structurepurple] {3^x}; \addlegendentry{$3^x$ (slope $>1$ at $0$)}
\addplot[dashed, gray, domain=-1.1:1.1] {1 + x};
\fill (axis cs:0,1) circle (2pt);
\node[anchor=north east] at (axis cs:0,1) {$(0,1)$};
\end{axis}
\end{tikzpicture}
\caption{Every exponential passes through $(0,1)$ but leaves at a different
slope. The base $e$ is the unique one whose tangent at $(0,1)$ is the line
$y = 1 + x$ (slope exactly $1$, dashed).}
\label{ch12:fig-exp-family}
\end{figure}

\subsection{Defining $e$ as the base that makes $m(b)=1$}

We have found the cleanest possible definition of the number $e$. Among all
exponential functions, exactly one has slope $1$ at $x=0$. Its base is what
we call $e$.

\begin{definition}[The number $e$, via slope]
\label{ch12:def-e}
$e$ is the unique base $b>0$ for which
\[
   \lim_{h\to 0}\frac{b^h-1}{h} = 1,
   \qquad\text{equivalently}\qquad
   \frac{\dd}{\dd x}b^x\Big|_{x=0} = 1.
\]
\end{definition}

With this base the constant $m(e)=1$, so the leftover limit disappears and
we land on the jewel of the chapter.

\begin{theorem}[The exponential is its own derivative]
\label{ch12:thm-ex}
\[
   \ddx{e^x} = e^x.
\]
\end{theorem}

\begin{proof}
From the general formula $\ddx{b^x}=m(b)\,b^x$ with $b=e$, and
$m(e)=1$ by Definition~\ref{ch12:def-e}, we get $\ddx{e^x}=1\cdot e^x = e^x$.
\end{proof}

\begin{happeningbox}
Look at what ``$\ddx{e^x}=e^x$'' really asserts. At every point $(x, e^x)$
on the curve, the steepness of the graph equals its height above the axis.
At height $1$ the slope is $1$; at height $7$ the slope is $7$; at height
$100$ the slope is $100$. The curve is racing upward exactly as fast as it
is tall. That self-referential loop -- ``my rate of growth is my current
size'' -- is what makes $e^x$ explode, and it is why the same function
appears whenever growth feeds on itself: interest on interest, infection
spreading person to person, a chain reaction. The symbol $e^x$ is just a
name for ``the function that is its own rate of change and equals $1$ at
$0$.'' Everything else about it follows.
\end{happeningbox}

\begin{historybox}
The story of exponentials and logarithms is the story of calculus liberating itself from polynomial geometry to conquer the physical universe. In the early seventeenth century, astronomical navigation and planetary calculations (such as Johannes Kepler's calculations of planetary orbits) were paralyzed by the sheer drudgery of multiplying and dividing eight-digit numbers by hand. In 1614, John Napier published \emph{Mirifici Logarithmorum Canonis Descriptio}, followed shortly by Henry Briggs's 1624 tables of base-10 logarithms in \emph{Arithmetica Logarithmica}. By converting arduous multiplications into simple additions ($\log(ab) = \log a + \log b$), logarithms were celebrated as having ``doubled the lives of astronomers.''

Yet for more than a century, logarithms were regarded purely as computational tables and geometric areas under hyperbolas, not as functions, and certainly not as the inverses of an exponential. The decisive transformation came in 1748 with Leonhard Euler's monumental two-volume treatise \emph{Introductio in analysin infinitorum}. Euler shifted the foundation of mathematics away from geometric curves and onto analytical \emph{functions}. He systematically established the number $e \approx 2.718281828459$, demonstrating that the compound-interest limit $\lim_{n\to\infty}(1 + 1/n)^n$, the infinite power series $\sum_{k=0}^\infty 1/k!$, and the unique base whose exponential slope is $1$ at the origin are all identical manifestations of the same constant.

This new analytical world was not without fierce controversy. A celebrated dispute erupted between Euler and his former mentor, Johann Bernoulli, over the logarithm of negative numbers. Bernoulli insisted that since $(-x)^2 = x^2$, taking logarithms must yield $2\ln(-x) = 2\ln(x)$, implying that $\ln(-x) = \ln(x)$ for all real $x$---so negative numbers would have ordinary, real logarithms! Euler dismantled Bernoulli's argument in 1749 by demonstrating that the logarithm is intrinsically multi-valued in the complex plane. Euler showed that $\ln(-1) = (2k+1)\pi i$ for any integer $k$; negative inputs rotate into the complex domain by an odd multiple of $\pi$ radians. The ``natural'' logarithm was far deeper than either Napier or Bernoulli had dreamed: it tied real rates of change directly to rotational geometry in the complex plane ($e^{ix} = \cos x + i\sin x$).
\end{historybox}

\begin{connectionbox}
Why does the slope-$1$ base equal the compound-interest limit? Here is the
quick reconciliation. Suppose $b$ is the base with $m(b)=1$, i.e.
$(b^h-1)/h \to 1$ as $h\to 0$. For small $h$ this says $b^h - 1 \approx h$,
so $b^h \approx 1+h$, hence $b \approx (1+h)^{1/h}$. Letting $h = 1/n$ and
$n\to\infty$ turns the right side into $\lim_{n\to\infty}(1+\tfrac1n)^n$ --
exactly the compound-interest definition of $e$. The two definitions are the
same fact viewed from two heights.
\end{connectionbox}

\section{The Natural Logarithm by Inverse Reasoning}
\label{ch12:ln}

Now that $e^x$ is tamed, its inverse comes almost for free. Recall that
$\ln x$ (the natural logarithm) is \emph{defined} as the inverse of $e^x$:
\[
   y = \ln x \quad\Longleftrightarrow\quad e^y = x,\qquad (x>0).
\]
So $\ln$ undoes $e^{(\cdot)}$ and vice versa: $e^{\ln x} = x$ and
$\ln(e^x)=x$. We differentiate it two ways -- once with a slick implicit
argument, and once straight from the limit -- because each illuminates the
other.

\subsection{Derivative of $\ln x$ by implicit differentiation}

\begin{intuitionbox}
A function and its inverse are reflections of each other across the line
$y=x$. Reflection swaps the roles of $x$ and $y$, and it swaps ``run'' with
``rise.'' So wherever $e^x$ has a steep slope, its mirror image $\ln x$ has
a \emph{shallow} one --- the slope of the inverse is the \emph{reciprocal} of
the slope of the original: $\big(f^{-1}\big)'(x) = 1 / f'\!\big(f^{-1}(x)\big)$.
Swapping input and output coordinates inverts the tangent slope directly;
the algebraic computation below confirms this geometric reflection.
\end{intuitionbox}

Start from the relation $e^y = x$ and differentiate both sides with respect
to $x$. The left side needs the chain rule, since $y$ is a function of $x$:
\[
   \ddx{e^y} = \ddx{x}
   \;\;\Longrightarrow\;\;
   e^y \cdot \dv{y}{x} = 1.
\]
Solve for $\dv{y}{x}$ and then replace $e^y$ by $x$ (that is what $e^y$
equals):
\[
   \dv{y}{x} = \frac{1}{e^y} = \frac{1}{x}.
\]
That is the result, clean and complete.

\begin{theorem}[Derivative of the natural log]
\label{ch12:thm-ln}
For $x>0$,
\[
   \ddx{\ln x} = \frac{1}{x}.
\]
\end{theorem}

\begin{whybox}
The inverse-function trick works because differentiating the
\emph{identity} $e^{\ln x}=x$ forces the answer. Apply the chain rule to the
left side: $\ddx{e^{\ln x}} = e^{\ln x}\cdot(\ln x)' = x\cdot(\ln x)'$. The
right side is just $1$. So $x\cdot(\ln x)' = 1$, giving $(\ln x)' = 1/x$
with no guessing at all. Whenever you know a function's derivative and want
its inverse's derivative, differentiate the ``they cancel'' identity and
solve. The reciprocal-of-the-slope picture and this one-line computation are
the same idea: $\big(f^{-1}\big)'(x) = 1/f'\!\big(f^{-1}(x)\big)$.
\end{whybox}

\begin{figure}[H]
\centering
\begin{tikzpicture}
\begin{axis}[
  width=9.5cm, height=7.5cm,
  axis lines=middle, axis equal image,
  xlabel=$x$, ylabel=$y$,
  xmin=-0.5, xmax=4.2, ymin=-1.6, ymax=4.2,
  domain=0.08:4, samples=120,
  legend pos=south east, legend cell align=left,
  every axis plot/.append style={very thick},
]
\addplot[tangentred, domain=-1.4:1.4] {exp(x)};
\addlegendentry{$e^x$}
\addplot[curveblue] {ln(x)};
\addlegendentry{$\ln x$}
\addplot[gray, dashed, domain=-0.5:4.2] {x};
\addlegendentry{$y=x$}
\end{axis}
\end{tikzpicture}
\caption{$\ln x$ is the mirror image of $e^x$ across $y=x$. Where $e^x$
climbs steeply, $\ln x$ rises gently: the slopes are reciprocals. As
$x\to 0^+$, $\ln x$ plunges to $-\infty$ ever more steeply, matching
$1/x\to+\infty$.}
\label{ch12:fig-ln-reflection}
\end{figure}

\subsection{The same result straight from the limit}

The same $1/x$ also drops out of the raw difference quotient,
with no inverse machinery -- and the derivation secretly re-derives the
compound-interest formula for $e$ along the way. We compute:
\begin{align*}
\ddx{\ln x}
&= \lim_{h\to 0}\frac{\ln(x+h)-\ln x}{h}
 = \lim_{h\to 0}\frac{1}{h}\ln\!\left(\frac{x+h}{x}\right)
 = \lim_{h\to 0}\frac{1}{h}\ln\!\left(1+\frac{h}{x}\right).
\end{align*}
Let $u = h/x$, so $h = ux$ and $u\to 0$ as $h\to 0$. Then $1/h = 1/(ux)$:
\[
\ddx{\ln x}
 = \lim_{u\to 0}\frac{1}{ux}\,\ln(1+u)
 = \frac{1}{x}\,\lim_{u\to 0}\frac{\ln(1+u)}{u}
 = \frac{1}{x}\,\lim_{u\to 0}\ln\!\Big[(1+u)^{1/u}\Big].
\]
Inside the bracket sits $(1+u)^{1/u}$, which tends to $e$ as $u\to 0$ (the
compound-interest limit again, with $u=1/n$). Since $\ln$ is continuous we
may pass the limit inside: the bracket becomes $\ln e = 1$. Therefore
\[
   \ddx{\ln x} = \frac{1}{x}\cdot \ln e = \frac{1}{x}.
\]
Both roads lead to $1/x$, and both quietly rely on the same defining
property of $e$. That is reassuring: the formula is not an artifact of one
clever trick.

\begin{aestheticbox}
There is a small mystery hiding here that will pay off enormously in the
integration chapters. The power rule gives $\ddx{x^n} = n x^{n-1}$ for every
exponent $n$ -- \emph{except} $n=0$, where it produces $0\cdot x^{-1}=0$
instead of $x^{-1}$. So in the entire family $\{x^n\}$ there is no function
whose derivative is $x^{-1}=1/x$. Nature abhors the gap, and fills it with a
function from a completely different family: $\ln x$. This is why
$\int \frac{1}{x}\,\dd x = \ln\abs{x}+C$ -- the logarithm is exactly the
``missing'' antiderivative that the power rule could never supply.
\end{aestheticbox}

\section{General Bases: $a^x$ and $\log_a x$}
\label{ch12:general-base}

We deliberately built everything on base $e$ because it is the base with no
clutter. Other bases inherit their derivatives from $e$ through one rewrite.

\subsection{Derivative of $a^x$}

The trick is to express \emph{any} exponential in terms of the natural one.
Since $a = e^{\ln a}$, we have $a^x = \big(e^{\ln a}\big)^x = e^{(\ln a)x}$.
Now it is just the chain rule with inner function $(\ln a)x$:
\[
   \ddx{a^x}
   = \ddx{e^{(\ln a)x}}
   = e^{(\ln a)x}\cdot \ddx{(\ln a)x}
   = a^x \cdot \ln a.
\]

\begin{focusbox}
\[
   \ddx{a^x} = a^x \ln a.
\]
Compare with our very first computation, $\ddx{a^x}=m(a)\,a^x$. We have now
\emph{identified} the mystery constant: $m(a) = \ln a$. The slope of $a^x$ at
the $y$-axis is the natural log of its base. So $m(2)=\ln 2\approx 0.6931$
and $m(3)=\ln 3\approx 1.0986$ -- right where the finite-difference estimates
from the lab box ($0.6934$ and $1.0992$) were pointing. The small gap is just
the overshoot from using $h=0.001$ instead of the true limit; shrink $h$ and
the estimates march straight onto $\ln 2$ and $\ln 3$.
\end{focusbox}

\begin{labbox}
Check the loop closes. Our finite-difference experiment gave
$m(2)\approx 0.6934$. Now the theory says $m(2)=\ln 2 = 0.693147\ldots$ --
spot on, the last digit's drift being only the $h=0.001$ overshoot. And
$m(e)=\ln e = 1$, recovering the defining property of $e$ exactly. The
numerical experiment and the clean theory agree as closely as you care to
push the step size down.
\end{labbox}

\subsection{Derivative of $\log_a x$}

The general logarithm reduces to the natural one by the change-of-base
formula $\log_a x = \dfrac{\ln x}{\ln a}$. Since $\ln a$ is a constant,
\[
   \ddx{\log_a x}
   = \frac{1}{\ln a}\,\ddx{\ln x}
   = \frac{1}{(\ln a)\,x}.
\]
Once more $e$ is the lucky base: with $a=e$ the factor $\ln a=1$ vanishes
and we recover $1/x$. Every stray $\ln a$ floating around in a calculus
formula is the universe nagging you to switch to base $e$.

\section{Logarithmic Differentiation}
\label{ch12:logdiff}

Here is a technique that turns the logarithm from an object of study into a
\emph{tool}. As an algebraic simplifier the natural log does two things:
it converts products into sums and exponents into multipliers,
\[
   \ln(fg) = \ln f + \ln g,
   \qquad
   \ln\!\big(f^{\,g}\big) = g\,\ln f.
\]
Differentiation loves sums and hates products and towers of exponents. So
the plan is: take $\ln$ of both sides \emph{first} to break the mess apart,
differentiate the easy pieces, then multiply back.

\begin{whybox}
The engine is the chain rule applied to $\ln y$. If $y=y(x)$, then
\[
   \ddx{\ln y} = \frac{1}{y}\,\dv{y}{x} = \frac{y'}{y}.
\]
The quantity $y'/y$ is the \emph{logarithmic derivative}, also known as the
\emph{relative} or \emph{fractional} rate of change -- growth measured as a
percentage of current size. So whenever you compute $\frac{d}{dx}\ln y$ you
are really computing $y'/y$; multiply by $y$ at the end to recover $y'$
itself. Logarithmic differentiation is just this identity used in reverse:
build $y'/y$ out of easy log-pieces, then scale up by $y$.
\end{whybox}

\begin{example}[A product of three factors]
\label{ch12:ex-logdiff-product}
Differentiate
\(
   y = \dfrac{(x^2+1)^3\,\sqrt{x-4}}{(x+2)^5}.
\)
Doing this with the quotient and product rules directly would be a
bookkeeping nightmare. Instead take logs:
\[
   \ln y = 3\ln(x^2+1) + \tfrac12\ln(x-4) - 5\ln(x+2).
\]
Differentiate term by term (each is an easy chain-rule piece):
\[
   \frac{y'}{y}
   = 3\cdot\frac{2x}{x^2+1} + \frac12\cdot\frac{1}{x-4} - 5\cdot\frac{1}{x+2}
   = \frac{6x}{x^2+1} + \frac{1}{2(x-4)} - \frac{5}{x+2}.
\]
Finally multiply back by $y$:
\[
   y' = \frac{(x^2+1)^3\sqrt{x-4}}{(x+2)^5}
        \left[\frac{6x}{x^2+1} + \frac{1}{2(x-4)} - \frac{5}{x+2}\right].
\]
All the structure of the original function survives out front, with a tidy
sum of corrections inside the brackets.
\end{example}

\begin{example}[A variable base \emph{and} a variable exponent]
\label{ch12:ex-tower}
Differentiate $y = x^x$ for $x>0$. Notice neither the power rule
($\ddx{x^n}=nx^{n-1}$, exponent constant) nor the exponential rule
($\ddx{a^x}=a^x\ln a$, base constant) applies -- here \emph{both} the base
and the exponent vary. Logs untangle it. Take $\ln$:
\[
   \ln y = x\ln x.
\]
Differentiate the right side with the product rule:
\[
   \frac{y'}{y} = 1\cdot \ln x + x\cdot\frac1x = \ln x + 1.
\]
Multiply back: $y' = x^x(\ln x + 1)$. As a sanity check, $x^x$ has its
minimum where $y'=0$, i.e. $\ln x = -1$, i.e. $x = 1/e \approx 0.368$ --
which is indeed where the graph of $x^x$ bottoms out
(Figure~\ref{ch12:fig-xtox}).
\end{example}

\begin{figure}[H]
\centering
\begin{tikzpicture}
\begin{axis}[
  width=10cm, height=6cm,
  axis lines=middle,
  xlabel=$x$, ylabel=$y$,
  xmin=0, xmax=1.7, ymin=0, ymax=1.7,
  domain=0.02:1.6, samples=120,
  legend pos=north west, legend cell align=left,
  every axis plot/.append style={very thick},
]
\addplot[curveblue] {exp(x*ln(x))}; \addlegendentry{$y=x^x$}
\addplot[only marks, mark=*, tangentred] coordinates {(0.36788,0.69220)};
\draw[dashed, gray] (axis cs:0.36788,0) -- (axis cs:0.36788,0.69220);
\node[anchor=north, font=\small] at (axis cs:0.36788,0) {$1/e$};
\end{axis}
\end{tikzpicture}
\caption{The graph of $x^x$ for $x>0$. Logarithmic differentiation gives
$y' = x^x(\ln x+1)$, which vanishes only at $x=1/e\approx 0.368$, the marked
low point. There $x^x = (1/e)^{1/e}\approx 0.692$.}
\label{ch12:fig-xtox}
\end{figure}

\begin{trapbox}
Do not write $\ddx{x^x} = x\cdot x^{x-1}$ (treating the exponent as
constant), and do not write $\ddx{x^x} = x^x\ln x$ (treating the base as
constant). \emph{Both} are wrong because \emph{both} the base and the
exponent move. The correct derivative $x^x(\ln x+1)$ contains a piece from
each: the $1$ comes from varying the base, the $\ln x$ from varying the
exponent. When in doubt with anything of the form
$(\text{stuff})^{(\text{stuff})}$, take logs.
\end{trapbox}

\section{Sine and Cosine: the Two Limits That Decide Everything}
\label{ch12:trig}

Now to the periodic functions. Their derivatives rest on two famous limits.
Everything else is the angle-addition formula and bookkeeping. The two
limits are
\[
   \lim_{\theta\to 0}\frac{\sin\theta}{\theta} = 1,
   \qquad
   \lim_{\theta\to 0}\frac{\cos\theta - 1}{\theta} = 0,
\]
with $\theta$ measured in \emph{radians}. We trust these from the limits
chapter, but the unit-circle picture makes them believable, so look
before we compute.

\begin{intuitionbox}
On the unit circle, the \emph{angle} $\theta$ (in radians) is the arc length
swept out, and $\sin\theta$ is the vertical height of the endpoint. For a
tiny angle, the arc and the vertical rise are nearly the same little
segment, so $\sin\theta \approx \theta$ -- their ratio heads to $1$.
Meanwhile $\cos\theta$ is the horizontal coordinate, which for a tiny angle
is barely less than $1$; it drops away only to \emph{second} order
($\cos\theta\approx 1-\theta^2/2$), so $(\cos\theta-1)/\theta\approx
-\theta/2\to 0$. Height tracks angle to first order; width clings to $1$.
\end{intuitionbox}

\begin{figure}[H]
\centering
\begin{tikzpicture}[scale=0.92]
  % Left Panel: Unit circle geometry
  \begin{scope}[shift={(0,0)}]
    \node[font=\footnotesize\scshape, text=gray!70!black] at (2.2, 4.75) {Unit Circle Coordinates};
    \draw[gray!40, thick] (2.0,1.8) circle (1.8);
    \draw[->, >=stealth, gray!60] (0.0,1.8) -- (4.2,1.8) node[right, font=\small] {$x$};
    \draw[->, >=stealth, gray!60] (2.0,0.0) -- (2.0,3.9) node[above, font=\small] {$y$};

    \coordinate (O) at (2.0,1.8);
    \coordinate (P) at ({2.0 + 1.8*cos(40)}, {1.8 + 1.8*sin(40)});

    \draw[curveblue, very thick] (O) -- (P);
    \draw[tangentred, very thick] ({2.0 + 1.8}, 1.8) arc (0:40:1.8);
    \node[tangentred, font=\small] at (3.2, 2.2) {$\theta$};

    \draw[structurepurple, thick, dashed] (P) -- ({2.0 + 1.8*cos(40)}, 1.8);
    \node[structurepurple, font=\small, right] at ({2.0 + 1.8*cos(40)}, {1.8 + 0.9*sin(40)}) {$\sin\theta$};

    \draw[targetgreen, very thick] (O) -- ({2.0 + 1.8*cos(40)}, 1.8);
    \node[targetgreen, font=\small, below] at ({2.0 + 0.9*cos(40)}, 1.8) {$\cos\theta$};

    \fill[black] (P) circle (2pt) node[above right, font=\small] {$(\cos\theta, \sin\theta)$};
  \end{scope}

  % Divider
  \draw[gray!30, dashed] (5.4, 0.0) -- (5.4, 4.9);

  % Right Panel: Velocity Vector & Derivative Components
  \begin{scope}[shift={(6.2,0)}]
    \node[font=\footnotesize\scshape, text=gray!70!black] at (3.5, 4.75) {Velocity Vector \& Differential Step};

    % Circle arc around P
    \draw[gray!40, thick] (1.0, 1.8) arc (15:70:3.2);
    \coordinate (Q) at (2.6, 2.5);
    \fill[black] (Q) circle (2pt) node[below right, font=\small] {$P(\cos\theta, \sin\theta)$};

    % Velocity vector (tangent to circle, rotated 90 deg counterclockwise)
    \coordinate (V) at ({2.6 + 1.8*cos(130)}, {2.5 + 1.8*sin(130)});
    \draw[->, >=stealth, very thick, tangentred] (Q) -- (V) node[above right=1pt, font=\footnotesize] {$\vec{v} = (-\sin\theta, \cos\theta)$};

    % Differential triangle from Q to V
    \coordinate (H) at ({2.6 + 1.8*cos(130)}, 2.5);
    \draw[targetgreen!80!black, thick, dashed] (Q) -- (H) node[midway, below, font=\small] {$-\sin\theta\,\dd\theta$};
    \draw[structurepurple, thick, dashed] (H) -- (V) node[midway, left, font=\small] {$\cos\theta\,\dd\theta$};

    % Explanation card
    \draw[sourceblue!70, fill=sourceblue!10, rounded corners=2pt] (0.2, 0.1) rectangle (6.8, 1.3);
    \node[anchor=west, font=\small] at (0.4, 0.95) {\textbf{Horizontal rate}: $\dv{}{\theta}(\cos\theta) = -\sin\theta$};
    \node[anchor=west, font=\small] at (0.4, 0.45) {\textbf{Vertical rate}: $\dv{}{\theta}(\sin\theta) = +\cos\theta$};
  \end{scope}
\end{tikzpicture}
\caption{The trigonometric derivatives visualized on the unit circle. Left: The coordinates of a point
at angle $\theta$ are $(\cos\theta, \sin\theta)$. Right: A counterclockwise nudge along the circle
of arc length $\dd\theta$ moves along the tangent velocity vector $\vec{v} = (-\sin\theta, \cos\theta)$, perpendicular
to the radius. Resolving this displacement into horizontal and vertical components reveals why
$\dd(\cos\theta)/\dd\theta = -\sin\theta$ and $\dd(\sin\theta)/\dd\theta = \cos\theta$.}
\label{ch12:fig-unit-circle}
\end{figure}

\subsection{Differentiating sine}

Now we run the difference quotient for $\sin x$ and use the angle-addition
formula $\sin(x+h) = \sin x\cos h + \cos x\sin h$.

\begin{align*}
\ddx{\sin x}
&= \lim_{h\to 0}\frac{\sin(x+h)-\sin x}{h}\\[2pt]
&= \lim_{h\to 0}\frac{\sin x\cos h + \cos x\sin h - \sin x}{h}\\[2pt]
&= \lim_{h\to 0}\left[\sin x\cdot\frac{\cos h - 1}{h}
                    + \cos x\cdot\frac{\sin h}{h}\right]\\[2pt]
&= \sin x\cdot\underbrace{\lim_{h\to0}\frac{\cos h-1}{h}}_{=\,0}
 + \cos x\cdot\underbrace{\lim_{h\to0}\frac{\sin h}{h}}_{=\,1}\\[2pt]
&= \sin x\cdot 0 + \cos x\cdot 1 \;=\; \cos x.
\end{align*}

\begin{theorem}[Derivatives of sine and cosine]
\label{ch12:thm-sincos}
\[
   \ddx{\sin x} = \cos x,
   \qquad
   \ddx{\cos x} = -\sin x.
\]
\end{theorem}

The cosine derivative falls out the same way, with
$\cos(x+h)=\cos x\cos h - \sin x\sin h$. Splitting and taking limits,
\[
\ddx{\cos x}
= \cos x\cdot\Big(\lim_{h\to0}\tfrac{\cos h-1}{h}\Big)
 - \sin x\cdot\Big(\lim_{h\to0}\tfrac{\sin h}{h}\Big)
= \cos x\cdot 0 - \sin x\cdot 1 = -\sin x.
\]

\begin{happeningbox}
The two limits did all the work, and each played a clear role. The limit
$\sin h/h\to 1$ is what makes the derivative have its full strength -- it is
the reason $\ddx{\sin x}$ is the \emph{whole} cosine and not some shrunken
multiple of it. That ``$1$'' is the trigonometric analogue of $m(e)=1$ for
the exponential: it is the payoff for measuring angles in radians. The other
limit, $(\cos h-1)/h\to 0$, says cosine is \emph{flat} at its peak -- it
contributes nothing to first order, which is why no stray $\sin x$ term
pollutes the answer. If you measured angles in degrees instead, $\sin h/h$
would tend to $\pi/180$, not $1$, and every trig derivative would drag an
ugly factor of $\pi/180$ forever. Radians are the natural unit for exactly
the same reason $e$ is the natural base: they make the leftover constant
equal to $1$.
\end{happeningbox}

\begin{trapbox}
The radian warning is not academic. If your calculator is in degree mode,
$\ddx{\sin x}\neq\cos x$ -- you would need $\frac{\pi}{180}\cos x$. All of
calculus assumes radians. A surprising number of wrong answers in physics
and engineering trace back to a calculator left in degree mode.
\end{trapbox}

\begin{figure}[H]
\centering
\begin{tikzpicture}
\begin{axis}[
  width=12cm, height=6cm,
  axis lines=middle,
  xlabel=$x$, ylabel=$y$,
  xmin=-0.4, xmax=6.6, ymin=-1.4, ymax=1.4,
  xtick={0,1.5708,3.1416,4.7124,6.2832},
  xticklabels={$0$,$\frac{\pi}{2}$,$\pi$,$\frac{3\pi}{2}$,$2\pi$},
  domain=0:6.2832, samples=160,
  legend pos=south west, legend cell align=left,
  every axis plot/.append style={very thick},
]
\addplot[curveblue] {sin(deg(x))};   \addlegendentry{$\sin x$}
\addplot[tangentred] {cos(deg(x))};  \addlegendentry{$\cos x = (\sin x)'$}
\end{axis}
\end{tikzpicture}
\caption{Where $\sin x$ rises most steeply (at $x=0$, slope $1$) its
derivative $\cos x$ peaks at $1$; where $\sin x$ flattens at its crest
($x=\pi/2$) the derivative crosses zero. The slope of $\sin$ \emph{is}
$\cos$, point by point.}
\label{ch12:fig-sin-cos}
\end{figure}

\begin{aestheticbox}
Differentiating sine and cosine sends you around a little cycle of length
four:
\[
   \sin x \xrightarrow{\;d\;} \cos x \xrightarrow{\;d\;} -\sin x
   \xrightarrow{\;d\;} -\cos x \xrightarrow{\;d\;} \sin x.
\]
Four derivatives bring you home. This is the seed of oscillation: anything
satisfying $y'' = -y$ must be built from sines and cosines, and that
equation governs every pendulum, spring, and wave in physics. Contrast the
exponential, whose derivative cycle has length \emph{one} ($e^x\to e^x$):
exponentials grow, trig functions circle. Later, Euler's formula
$e^{i\theta}=\cos\theta + i\sin\theta$ will reveal these as the \emph{same}
phenomenon viewed along the real and imaginary axes -- circling is just
exponential growth in an imaginary direction.
\end{aestheticbox}

\subsection{The rest of the trig functions}

With $\sin$ and $\cos$ in hand, the other four come from the quotient rule
and the Pythagorean identity. Take the tangent:
\[
   \ddx{\tan x}
   = \ddx{\frac{\sin x}{\cos x}}
   = \frac{(\cos x)(\cos x) - (\sin x)(-\sin x)}{\cos^2 x}
   = \frac{\cos^2 x + \sin^2 x}{\cos^2 x}
   = \frac{1}{\cos^2 x} = \sec^2 x.
\]
The numerator collapsed to $1$ because of $\sin^2+\cos^2=1$ -- a small
miracle of cancellation. The same quotient-rule routine gives the rest.

\begin{focusbox}
\[
\begin{array}{ll}
\ddx{\sin x} = \cos x, &\quad \ddx{\cos x} = -\sin x, \\[6pt]
\ddx{\tan x} = \sec^2 x, &\quad \ddx{\cot x} = -\csc^2 x, \\[6pt]
\ddx{\sec x} = \sec x\tan x, &\quad \ddx{\csc x} = -\csc x\cot x.
\end{array}
\]
Notice the pattern: every ``co-'' function ($\cos,\cot,\csc$) carries a minus
sign in its derivative. That bookkeeping minus is a reliable memory aid.
\end{focusbox}

\begin{example}[Using the chain rule with trig]
\label{ch12:ex-chain-trig}
Differentiate $g(x) = \sin(3x^2+1)$ and $h(x)=\sec(5x)$.
\[
   g'(x) = \cos(3x^2+1)\cdot \ddx{3x^2+1} = 6x\cos(3x^2+1),
\]
\[
   h'(x) = \sec(5x)\tan(5x)\cdot 5 = 5\sec(5x)\tan(5x).
\]
The trig derivative supplies the ``outer'' factor; the chain rule supplies
the derivative of the inside.
\end{example}

\section{Derivatives of the Inverse Trig Functions}
\label{ch12:inverse-trig}

The inverse trig functions $\arcsin$, $\arccos$, $\arctan$ answer ``which
angle has this sine/cosine/tangent?'' We differentiate them exactly as we
did $\ln x$: implicitly, using the cancellation identity, plus one
Pythagorean simplification at the end. The recurring surprise is that these
derivatives are \emph{algebraic} -- no trig left in the answer at all.

\subsection{Arcsine}

Let $y = \arcsin x$, so $\sin y = x$ with $y\in[-\tfrac\pi2,\tfrac\pi2]$.
Differentiate $\sin y = x$ implicitly:
\[
   \cos y\cdot\dv{y}{x} = 1
   \quad\Longrightarrow\quad
   \dv{y}{x} = \frac{1}{\cos y}.
\]
Now convert $\cos y$ into something in terms of $x$. From
$\sin^2 y + \cos^2 y = 1$ and $\sin y = x$, we get
$\cos y = \sqrt{1-x^2}$. (We take the $+$ root because on
$[-\tfrac\pi2,\tfrac\pi2]$ cosine is nonnegative -- the restricted range is
doing real work here.) Therefore
\[
   \ddx{\arcsin x} = \frac{1}{\sqrt{1-x^2}},\qquad -1<x<1.
\]

\begin{whybox}
Where did the square root come from? The implicit step gives a clean
$1/\cos y$, but the answer should be a function of $x$, not of the hidden
angle $y$. The right triangle with opposite side $x$ and hypotenuse $1$ has
$\sin y = x$, so by Pythagoras its adjacent side is $\sqrt{1-x^2}$, and that
adjacent-over-hypotenuse is $\cos y$. Drawing the triangle is the direct way
to turn ``$\cos(\arcsin x)$'' into ``$\sqrt{1-x^2}$.'' Every inverse-trig
derivative comes out algebraic for this reason: a trig function of an
inverse-trig function is always just a side ratio in a reference triangle.
\end{whybox}

\begin{figure}[H]
\centering
\begin{tikzpicture}[scale=0.92]
  % Panel 1: Arcsine reference triangle
  \begin{scope}[shift={(0,0)}]
    \node[font=\footnotesize\scshape, text=sourceblue!80!black] at (3.2, 3.8) {Arcsine: $y = \arcsin x \implies \sin y = x$};

    \coordinate (A) at (0.5, 0.5);
    \coordinate (B) at (4.8, 0.5);
    \coordinate (C) at (4.8, 3.0);

    \draw[thick, gray!70] (A) -- (B) -- (C) -- cycle;
    \draw[gray!60] (4.4, 0.5) -- (4.4, 0.9) -- (4.8, 0.9); % right angle

    % Sides
    \draw[tangentred, very thick] (A) -- (C) node[midway, above left, font=\small] {$\text{hypotenuse} = 1$};
    \draw[targetgreen!80!black, very thick] (B) -- (C) node[midway, right, font=\small] {$\text{opp} = x$};
    \draw[sourceblue, very thick] (A) -- (B) node[midway, below, font=\small] {$\text{adj} = \sqrt{1 - x^2}$};

    % Angle y
    \draw[quotientorange, thick] (1.3, 0.5) arc (0:30:0.8);
    \node[quotientorange!80!black, font=\small] at (1.6, 0.75) {$y$};

    % Deduction box
    \draw[sourceblue!70, fill=sourceblue!8, rounded corners=2pt] (0.2, -1.2) rectangle (5.8, -0.2);
    \node[font=\small] at (3.0, -0.7) {$\cos y = \sqrt{1-x^2} \implies \dv{y}{x} = \dfrac{1}{\cos y} = \dfrac{1}{\sqrt{1-x^2}}$};
  \end{scope}

  % Divider
  \draw[gray!30, dashed] (6.4, -1.4) -- (6.4, 4.0);

  % Panel 2: Arctangent reference triangle
  \begin{scope}[shift={(7.0,0)}]
    \node[font=\footnotesize\scshape, text=targetgreen!80!black] at (3.2, 3.8) {Arctangent: $y = \arctan x \implies \tan y = x$};

    \coordinate (D) at (0.5, 0.5);
    \coordinate (E) at (4.4, 0.5);
    \coordinate (F) at (4.4, 3.0);

    \draw[thick, gray!70] (D) -- (E) -- (F) -- cycle;
    \draw[gray!60] (4.0, 0.5) -- (4.0, 0.9) -- (4.4, 0.9); % right angle

    % Sides
    \draw[sourceblue, very thick] (D) -- (E) node[midway, below, font=\small] {$\text{adj} = 1$};
    \draw[targetgreen!80!black, very thick] (E) -- (F) node[midway, right, font=\small] {$\text{opp} = x$};
    \draw[tangentred, very thick] (D) -- (F) node[midway, above left, font=\small] {$\text{hyp} = \sqrt{1 + x^2}$};

    % Angle y
    \draw[quotientorange, thick] (1.3, 0.5) arc (0:33:0.8);
    \node[quotientorange!80!black, font=\small] at (1.6, 0.8) {$y$};

    % Deduction box
    \draw[targetgreen!70, fill=targetgreen!8, rounded corners=2pt] (0.2, -1.2) rectangle (5.8, -0.2);
    \node[font=\small] at (3.0, -0.7) {$\sec^2 y = 1+x^2 \implies \dv{y}{x} = \dfrac{1}{\sec^2 y} = \dfrac{1}{1+x^2}$};
  \end{scope}
\end{tikzpicture}
\caption{The reference triangles converting trig functions of inverse trig functions into pure algebra.
Left: For $y = \arcsin x$, a right triangle with hypotenuse $1$ and opposite side $x$ forces the adjacent
side to be $\sqrt{1-x^2}$ by the Pythagorean theorem, giving $\cos y = \sqrt{1-x^2}$.
Right: For $y = \arctan x$, adjacent side $1$ and opposite side $x$ force the hypotenuse to be $\sqrt{1+x^2}$,
giving $\sec^2 y = 1 + x^2$.}
\label{ch12:fig-invtrig-triangles}
\end{figure}

\subsection{Arctangent}

Let $y=\arctan x$, so $\tan y = x$. Differentiate:
\[
   \sec^2 y\cdot\dv{y}{x} = 1
   \quad\Longrightarrow\quad
   \dv{y}{x} = \frac{1}{\sec^2 y}.
\]
Use the identity $\sec^2 y = 1 + \tan^2 y = 1 + x^2$. So
\[
   \ddx{\arctan x} = \frac{1}{1+x^2},\qquad x\in\R.
\]
This one is a favorite because it is defined for all $x$, is beautifully
simple, and -- as you will see in the integration chapters -- it is exactly
the antiderivative needed to integrate $1/(1+x^2)$, the function whose area
hides $\pi$.

\begin{figure}[H]
\centering
\begin{tikzpicture}
\begin{axis}[
  width=11cm, height=6cm,
  axis lines=middle,
  xlabel=$x$, ylabel=$y$,
  xmin=-6, xmax=6, ymin=-2.1, ymax=2.2, clip=false,
  domain=-6:6, samples=160,
  legend pos=south east, legend cell align=left,
  every axis plot/.append style={very thick},
]
\addplot[curveblue] {rad(atan(x))}; \addlegendentry{$\arctan x$}
\addplot[tangentred] {1/(1+x^2)};   \addlegendentry{$\dfrac{1}{1+x^2}=(\arctan x)'$}
\addplot[gray, dashed, domain=-6:6] {1.5708};
\addplot[gray, dashed, domain=-6:6] {-1.5708};
\node[gray, anchor=south east] at (axis cs:6,1.5708) {$\tfrac\pi2$};
\end{axis}
\end{tikzpicture}
\caption{$\arctan x$ flattens toward the horizontal asymptotes
$\pm\pi/2$, and exactly where it flattens its derivative $1/(1+x^2)$ dies
away to $0$. The derivative peaks at $x=0$, where $\arctan$ is steepest.}
\label{ch12:fig-arctan}
\end{figure}

\begin{focusbox}
The inverse-trig derivative table (each on its natural domain):
\[
\begin{array}{ll}
\ddx{\arcsin x} = \dfrac{1}{\sqrt{1-x^2}}, &\quad
\ddx{\arccos x} = -\dfrac{1}{\sqrt{1-x^2}}, \\[10pt]
\ddx{\arctan x} = \dfrac{1}{1+x^2}, &\quad
\ddx{\operatorname{arccot} x} = -\dfrac{1}{1+x^2}, \\[10pt]
\ddx{\operatorname{arcsec} x} = \dfrac{1}{\abs{x}\sqrt{x^2-1}}, &\quad
\ddx{\operatorname{arccsc} x} = -\dfrac{1}{\abs{x}\sqrt{x^2-1}}.
\end{array}
\]
Again the ``co-'' partners differ only by a sign. That sign comes from
$\arccos x = \tfrac\pi2 - \arcsin x$: complementary angles add to a constant,
so their derivatives are negatives.
\end{focusbox}

\begin{connectionbox}
These derivatives matter most read \emph{backwards}, as integration
formulas. Differentiating the inverse trig functions is how we discover that
\[
   \int \frac{\dd x}{\sqrt{1-x^2}} = \arcsin x + C,
   \qquad
   \int \frac{\dd x}{1+x^2} = \arctan x + C.
\]
Two perfectly innocent-looking algebraic integrands turn out to have
\emph{trigonometric} antiderivatives. That is the deep reason
trigonometric substitution works for integrals containing $\sqrt{1-x^2}$ or
$1+x^2$ -- a topic we reach in the integration part. The angle is hiding
inside the algebra.
\end{connectionbox}

\section{A Gallery of Worked Examples}
\label{ch12:gallery}

To consolidate, here are the new rules mixed with the product, quotient,
and chain rules from the previous chapter.

\begin{example}[Exponential times trig]
\label{ch12:ex-damped}
Differentiate $f(x) = e^{-2x}\cos(3x)$, the shape of a damped oscillation.
Product rule, with each factor needing the chain rule:
\[
   f'(x) = \underbrace{(-2)e^{-2x}}_{\ddx{e^{-2x}}}\cos(3x)
         + e^{-2x}\underbrace{(-3\sin(3x))}_{\ddx{\cos(3x)}}
         = -e^{-2x}\big(2\cos 3x + 3\sin 3x\big).
\]
\end{example}

\begin{example}[Log of a composite]
\label{ch12:ex-log-comp}
Differentiate $f(x) = \ln\!\big(\sqrt{x^2+1}\big)$. Simplify with log laws
\emph{before} differentiating: $\ln\sqrt{x^2+1} = \tfrac12\ln(x^2+1)$. Then
\[
   f'(x) = \frac12\cdot\frac{2x}{x^2+1} = \frac{x}{x^2+1}.
\]
Always look for a log law to flatten the function first; it can turn a
chain-rule slog into a one-liner.
\end{example}

\begin{example}[Putting the base change to work]
\label{ch12:ex-base-change}
Differentiate $f(x) = 5^{\sin x}$. Rewrite as $e^{(\ln 5)\sin x}$, or just
apply $\ddx{a^u} = a^u\ln a\cdot u'$ directly with $a=5$, $u=\sin x$:
\[
   f'(x) = 5^{\sin x}\,(\ln 5)\,\cos x.
\]
\end{example}

\begin{example}[Inverse trig inside a chain]
\label{ch12:ex-arctan-chain}
Differentiate $f(x) = \arctan(e^x)$. Outer derivative $1/(1+u^2)$ with
$u=e^x$, times $u' = e^x$:
\[
   f'(x) = \frac{1}{1+(e^x)^2}\cdot e^x = \frac{e^x}{1+e^{2x}}.
\]
\end{example}

\begin{example}[Logarithmic differentiation with a root tower]
\label{ch12:ex-logdiff-2}
Differentiate $y = \dfrac{x^{3}\,e^{2x}}{\sqrt{1+x^2}}$. Take logs:
\[
   \ln y = 3\ln x + 2x - \tfrac12\ln(1+x^2).
\]
Differentiate:
\[
   \frac{y'}{y} = \frac{3}{x} + 2 - \frac{x}{1+x^2}.
\]
Multiply back:
\[
   y' = \frac{x^{3}e^{2x}}{\sqrt{1+x^2}}
        \left(\frac{3}{x} + 2 - \frac{x}{1+x^2}\right).
\]
\end{example}

\begin{walkbox}
Here is something to mull on a walk. We found three functions that are, in a
sense, ``eigenfunctions'' of differentiation -- the operation
$\ddx{}$ returns a simple multiple or relative of the input:
$e^x\to e^x$ (returns itself), and the pair $\sin,\cos$ that rotate among
themselves. Is there a function whose derivative is the function shifted, or
scaled by $x$, in some equally clean way? What kind of equation would
``$y' = $ (something simple involving $y$)'' have to be for its solution to
be elementary? You are circling the idea of a \emph{differential equation},
the subject the whole book builds toward. Every function in this chapter is
the answer to one: $y'=y$ gives $e^x$, $y''=-y$ gives $\sin$ and $\cos$,
$y'=1/x$ gives $\ln$.
\end{walkbox}

\begin{applicationbox}
These derivatives are the vocabulary of dynamics.
\begin{itemize}
\item \textbf{Radioactive decay / drug clearance.} If $N(t)=N_0 e^{-\lambda t}$,
then $N'(t) = -\lambda N(t)$: the decay rate is proportional to how much is
left. The derivative \emph{is} the law of decay.
\item \textbf{Exponential growth.} $\frac{\dd}{\dd t}(e^{kt})=ke^{kt}$ is why unchecked
populations, viral spread, and compound interest all share the same runaway
curve.
\item \textbf{Oscillation.} A mass on a spring obeys $x''=-\omega^2 x$, solved
by $x(t)=A\cos(\omega t)+B\sin(\omega t)$ precisely because of the
derivative cycle of sine and cosine.
\item \textbf{Machine learning.} The logistic ``sigmoid''
$\sigma(x)=1/(1+e^{-x})$ has the slick derivative
$\sigma'(x)=\sigma(x)\big(1-\sigma(x)\big)$ (try it!), which is why it is so
cheap to train neural networks with -- the gradient is built from values you
already computed.
\end{itemize}
\end{applicationbox}

\section{Summary}
\label{ch12:summary}

Every derivative in this chapter was forced by the \emph{definition} of the
function, not pulled from a hat.
\begin{itemize}
\item $\ddx{b^x} = (\ln b)\,b^x$, and $\ln b = \lim_{h\to0}(b^h-1)/h$ is the
slope of $b^x$ at the $y$-axis. The base $e$ is defined as the one making
that slope $1$, which is exactly why $\ddx{e^x}=e^x$.
\item $\ddx{\ln x}=1/x$ and $\ddx{\log_a x} = 1/(x\ln a)$, by inverse
reasoning or straight from the limit; $\ln x$ fills the gap the power rule
leaves at exponent $-1$.
\item Logarithmic differentiation turns products, quotients, and
variable-exponent towers into sums via $y'/y = \ddx{\ln y}$.
\item $\ddx{\sin x}=\cos x$ and $\ddx{\cos x}=-\sin x$, driven by
$\sin h/h\to1$ and $(\cos h-1)/h\to 0$ from the unit circle -- and radians
are what make those limits clean.
\item Inverse trig derivatives are \emph{algebraic}:
$(\arcsin x)'=1/\sqrt{1-x^2}$, $(\arctan x)'=1/(1+x^2)$, found by implicit
differentiation plus a Pythagorean triangle.
\end{itemize}
The recurring theme: $e$ and radians are the ``natural'' choices precisely
because they make a leftover constant equal to $1$. Choose them, and
calculus stops nagging you with stray factors.

\subsection{The Scope and Horizon of Transcendental Derivatives}
\label{ch12:what-unlocked}

\begin{connectionbox}
With exponentials, logarithms, and trigonometric functions brought into differential calculus, applied mathematics exploded in scope:
\begin{itemize}
  \item \emph{Natural growth and decay}: Any process where rate is proportional to size ($\dd y/\dd t = ky$) is solved in closed form by $y(t) = y_0 e^{kt}$, modeling radioactive decay, compound capital, and population dynamics.
  \item \emph{Oscillation and wave mechanics}: The derivative cycle $(\sin x)' = \cos x$ and $(\cos x)' = -\sin x$ explains why simple harmonic motion satisfies the second-order acceleration equation $\dd^2 y/\dd t^2 = -\omega^2 y$.
  \item \emph{Algebraic antiderivatives from geometry}: The remarkable discovery that inverse trigonometric derivatives are purely algebraic ($(\arctan x)' = 1/(1+x^2)$) provides the vital bridge allowing us to integrate rational expressions that would otherwise be insoluble.
\end{itemize}

Yet transcendental functions break purely algebraic solvability. In polynomial calculus, finding where a derivative vanishes or locating where two curves intersect requires solving polynomial equations, for which algebraic methods (factoring, the quadratic formula, Cardano's formula) exist. But the moment exponentials or trig functions enter, optimization and root-finding produce \emph{transcendental equations}:
\[
  x\,e^x = 1, \qquad x = \cos x, \qquad e^{-x^2} = x.
\]
None of these equations can be solved for $x$ using a finite combination of elementary algebraic operations. We can write down the exact derivative formula in seconds, yet algebra cannot tell us the numerical coordinate where that rate equals zero.

Because algebra hits a firm wall when solving transcendental equations, calculus must build algorithmic tools. In Chapter~19 (\emph{Newton's Method and Numerical Root-Finding}), we transform the derivative itself into an iterative engine that converges quadratically to the roots of transcendental equations, producing twenty decimal digits of accuracy in a handful of mechanical steps.
\end{connectionbox}

\section{Exercises}
\label{ch12:exercises}

\begin{exercise}[Warm-up derivatives]
\label{ch12:exr-warmup}
Differentiate each function.
\begin{enumerate}[label=(\alph*)]
\item $f(x) = 7e^x - 4\ln x$
\item $g(x) = e^{3x}\sin x$
\item $h(x) = \dfrac{\cos x}{x^2}$
\item $p(x) = \tan(2x) + \sec(2x)$
\item $q(x) = 3^x + x^3$ (mind which rule each term needs)
\item $r(x) = \ln(\cos x)$
\end{enumerate}
\end{exercise}

\begin{exercise}[Chain-rule combinations]
\label{ch12:exr-chain}
Differentiate.
\begin{enumerate}[label=(\alph*)]
\item $\sin(e^x)$
\item $e^{\sin x}$
\item $\arctan(3x)$
\item $\arcsin(x^2)$
\item $\ln\!\big(\ln x\big)$
\item $\cos^2(5x)$
\end{enumerate}
\end{exercise}

\begin{exercise}[Logarithmic differentiation]
\label{ch12:exr-logdiff}
Use logarithmic differentiation.
\begin{enumerate}[label=(\alph*)]
\item $y = (\sin x)^x$ for $0<x<\pi$.
\item $y = x^{\ln x}$.
\item $y = \dfrac{(x+1)^4 (x-2)^3}{(x^2+5)^2}$.
\end{enumerate}
\end{exercise}

\begin{exercise}[Find the base]
\label{ch12:exr-base}
Numerically estimate $m(4)=\lim_{h\to0}(4^h-1)/h$ with $h=0.0001$, then
confirm it equals $\ln 4$ to three decimals. Explain in one sentence why the
estimate and $\ln 4$ should agree.
\end{exercise}

\begin{exercise}[The sigmoid identity]
\label{ch12:exr-sigmoid}
Let $\sigma(x) = \dfrac{1}{1+e^{-x}}$. Show that
$\sigma'(x) = \sigma(x)\big(1-\sigma(x)\big)$. (Hint: write
$\sigma = (1+e^{-x})^{-1}$ and use the chain rule, then refactor.)
\end{exercise}

\begin{exercise}[Tangent line]
\label{ch12:exr-tangent}
Find the equation of the tangent line to $y=\ln x$ at $x=e$. Where does this
tangent line cross the $x$-axis?
\end{exercise}

\begin{exercise}[A radian sanity check]
\label{ch12:exr-radian}
Suppose a student insists on degrees and defines $S(x)=\sin(x^\circ)$, i.e.
the sine of $x$ degrees. Using $x^\circ = \tfrac{\pi}{180}x$ radians and the
chain rule, show that $S'(x) = \tfrac{\pi}{180}\cos(x^\circ)$. Explain why
this confirms that radians are the natural unit for calculus.
\end{exercise}

\begin{exercise}[Conceptual: why proportional growth is exponential]
\label{ch12:exr-proportional}
Suppose a function $y(t)$ satisfies $y'(t) = k\,y(t)$ for a constant $k$, and
$y(0)=y_0$. Verify by differentiation that $y(t) = y_0 e^{kt}$ works. Then
argue informally why \emph{no} function of the form $y_0 b^t$ with
$b\ne e^k$ can satisfy the equation. (This is the precise sense in which
``rate proportional to size'' means ``exponential.'')
\end{exercise}

\begin{exercise}[Proof: derivative of $a^x$ from scratch]
\label{ch12:exr-proof-ax}
Without quoting $\ddx{a^x}=a^x\ln a$, prove it. Start from the difference
quotient, factor out $a^x$, and identify the leftover limit $m(a)$. Then,
using $a = e^{\ln a}$ and $\ddx{e^x}=e^x$ together with the chain rule, show
$m(a)=\ln a$. Comment on how the two methods give the same constant.
\end{exercise}

\begin{exercise}[Proof: arccosine and the sign]
\label{ch12:exr-proof-arccos}
Using $\arccos x = \tfrac{\pi}{2} - \arcsin x$, prove that
$\ddx{\arccos x} = -\dfrac{1}{\sqrt{1-x^2}}$. Then re-derive the same result
\emph{independently} by implicit differentiation of $\cos y = x$, being
careful about the sign of $\sin y$ on the range $y\in[0,\pi]$. Confirm the
two agree.
\end{exercise}

\begin{exercise}[Higher derivatives and the cycle]
\label{ch12:exr-higher}
Compute the first four derivatives of $f(x)=\sin x$ and confirm the
period-four cycle. Then find a simple formula for the $n$-th derivative
$f^{(n)}(x)$ in terms of $n \bmod 4$. What is $f^{(101)}(x)$?
\end{exercise}

\begin{exercise}[Mixed challenge]
\label{ch12:exr-challenge}
Differentiate $y = x^{\sin x}$ and find all $x$ in $(0,\pi)$ where the
relative rate of change $y'/y$ equals zero. Interpret these points on the
graph of $\ln y = (\sin x)(\ln x)$.
\end{exercise}
